\section{First contacts near an aligned endpoint}\label{sec:posteriors}

Proposition~\ref{prop:contact-upper} bounds the number of scales at which
the opposed tapes meet.  To obtain a competing lower bound, we consider a
positive bridge and a relative negative path sampled independently of
it.  We place the negative path one level below the bridge's endpoint.
The bridge's final scripted departure is exactly the
negative starting site, so there is always a contact one level below
the endpoint.  We shall show that,
after a prescribed script at a lower height, their first contact is
unlikely to be delayed until much higher up the bridge.

The lateral starting point of the negative path is determined by the
positive endpoint.  We keep this dependence in the conditional
probabilities of the possible endpoints, and estimate their running
maxima together.  The martingale and endpoint-posterior arguments below
adapt \cite[Lemma~6.1 and Proposition~6.2]{OpenAI2026}.  The following
lemma gives the needed bound.

\begin{lemma}[Maxima of a finite martingale family]\label{lem:posterior-max}
Let \(m\ge1\), and let \((w_i(e))_{i\ge0}\), \(1\le e\le m\), be
nonnegative martingales for the same filtration \((\mathcal F_i)_{i\ge0}\), with
\(\sum_{e=1}^m w_i(e)\le1\) almost surely for every \(i\).  Set
\[
 A_i(e)=\max_{0\le j\le i}w_j(e),\qquad
 A_i=\sum_{e=1}^m A_i(e),\qquad
 A_\infty=\lim_{i\to\infty}A_i,
 \qquad B=\log(m+1).
\]
Write also \(A_\infty(e)=\lim_i A_i(e)\).
There are absolute constants \(c,C>0\) such that
\begin{equation}\label{eq:posterior-exponential}
 \EE\exp(cA_\infty/B)\le C.
\end{equation}
Moreover, on any joint probability space with this martingale marginal,
every event \(V\) of probability \(p\) satisfies
\begin{equation}\label{eq:event-weighted-max}
 \EE[A_\infty\one_V]
 \le CBp\log(\mathrm e/p),
\end{equation}
where the right-hand side is defined to be zero at \(p=0\).
The event \(V\) need not belong to the martingale filtration.
\end{lemma}

The logarithmic first-moment bound underlying this lemma follows by
coordinatewise integration of the scalar maximal inequality; the same
posterior calculation appears in
\cite[Section~1.1]{KesselheimMolinaroPattonSingla2026}.
The passage from bounded conditional future increases to exponential
integrability is related to the energy inequalities for increasing
processes in \cite[Theorem~4 and its corollary]{Kikuchi1992}.
We give a direct threshold proof.

\begin{proof}
Fix a time \(i\).  Conditional on the filtration at that time, the
nonnegative-martingale maximal inequality
\cite[Chapter~V, first section, \S2, Theorem~1]{Ville1939}
bounds the probability that
coordinate \(e\) subsequently exceeds \(u>0\) by \(w_i(e)/u\).
For completeness, on a finite horizon this follows by stopping at the
first crossing, using the martingale identity and nonnegativity on the
complement of the crossing event; increasing the horizon gives the
infinite-horizon bound.  Since each coordinate is at most one,
integration over \(u\in[A_i(e),1]\) gives
\[
 \EE[A_\infty(e)-A_i(e)\mid\mathcal F_i]
 \le w_i(e)\log\frac1{A_i(e)}
 \le w_i(e)\log\frac1{w_i(e)}.
\]
If \(w_i(e)=0\), its future values vanish almost surely conditionally,
and the expected increase is zero.  Add the missing mass
\(1-\sum_e w_i(e)\) to form a probability vector.  Its entropy is at
most \(\log(m+1)\), so
\begin{equation}\label{eq:maximum-remaining-increase}
 \EE[A_\infty-A_i\mid\mathcal F_i]\le B.
\end{equation}

We also have \(A_0\le1\) and
\(A_i-A_{i-1}\le\sum_e w_i(e)\le1\).
Consider the thresholds \(b_j=1+j(2B+1)\), \(j\ge1\).
The probability of strictly crossing the first threshold is at most
\(1/2\), by \eqref{eq:maximum-remaining-increase} and Markov's
inequality.  At the first crossing of any threshold, the overshoot is
at most one.  Crossing the next threshold therefore requires a further
increase greater than \(2B\), whose conditional probability is at most
\(1/2\).  This argument at a crossing time follows by splitting over
its possible finite values.  It gives a geometric tail for
\(A_\infty\) at the thresholds \(b_j\), and hence
\eqref{eq:posterior-exponential}.

For \(p>0\), conditional Jensen gives
\[
 \exp\!\left(\frac cB\EE[A_\infty\mid V]\right)
 \le \EE[\exp(cA_\infty/B)\mid V]\le C/p.
\]
Taking logarithms and multiplying by \(p\) proves
\eqref{eq:event-weighted-max}, after changing the absolute constant.
\end{proof}

We now formulate the first-contact estimate.  For a path in absolute
coordinates, write \(T_j\) for its first hit of the hyperplane
\(\{x:x_1=j\}\), specifying the path when needed.  Recall that
\(\mathcal B_N^+\) is the positive exact-cut bridge law: its
concatenated path has the raw marked law through \(T_N\), conditioned
on \(\mathcal D_N^+\).  The event \(\mathcal D_k^+\) prescribes the
\(M\)-step script immediately after the first hit of \(k-M\), reached
before height \(-1\).

\begin{proposition}[First contact above a scripted prefix]
\label{prop:first-contact}
Assume \(p_+p_->0\), and let \(M\le k<r<N\) be integers with
\(r-k\ge R\).  Start a positive bridge \(Y\) with law
\(\mathcal B_N^+\) at the origin, and put
\(E_{\mathrm{lat}}=\pi Y_{T_N}\).
Independently of \(Y\), sample a relative marked path with law
\(\widehat P_-\), and translate it to start at
\((N-1,E_{\mathrm{lat}})\); denote the resulting path by \(\beta\).
Define
\[
 \tau=\inf\{i:T_k(Y)\le i<T_N(Y),\quad
       \|Y_i-\beta_j\|_\infty\le R\text{ for some }j\ge0\},
\]
with \(\inf\varnothing=\infty\).  Let
\begin{equation}\label{eq:first-contact-event}
 \mathcal C_{k,r,N}=
 \left\{
 \begin{gathered}
  \mathcal D_k^+(Y),\quad \tau<T_N(Y),\quad (Y_\tau)_1\ge r,\\
  \min_{T_k(Y)\le i\le\tau}(Y_i)_1\ge k,
  \quad |E_{\mathrm{lat}}|\le N^2
 \end{gathered}
 \right\}.
\end{equation}
Here \(\mathcal D_k^+(Y)\) is a condition on the prefix through the
first hit of \(k\).  There is a constant \(C\), independent of
\(k,r,N\), such that, with
\(\Delta=u_{r-k}^+-u_{N-k}^+\),
\begin{equation}\label{eq:first-contact-bound}
 P(\mathcal C_{k,r,N})
 \le C\log(2N)\,\Delta\log^2(\mathrm e/\Delta).
\end{equation}
The right-hand side is defined to be zero when \(\Delta=0\).
\end{proposition}

The small parameter in this estimate is controlled by the width tail:
Lemma~\ref{lem:width-tail} gives
\[
 \Delta\le C(N-r)F_+(r-k).
\]
Thus the remaining height \(N-r\) is compared with the width tail at
the progress \(r-k\) already made above the scripted prefix.

\begin{proof}
We first express the event using finite positive prefixes and the full
negative path.  This permits a comparison in a shared environment
without treating \(\tau\) as a stopping time.  The endpoint
conditional probabilities remain attached to the positive prefixes;
Lemma~\ref{lem:posterior-max} will control their contribution.

\emph{Prefix probabilities.}
Let
\[
 \mathcal E_N=\{e\in\ZZ^{d-1}:|e|\le N^2\}.
\]
Enumerate the marked prefixes \(a_0\) that realize
\(\mathcal D_k^+\), ending at their first hit \(z\) of height \(k\).
For each such \(a_0\), enumerate all finite marked words \(a\) from
\(z\) whose vertices have heights in \([k,N-1]\) and whose terminal
site \(y=y(a)\) has height at least \(r\).
For a marked prefix \(c\) and \(e\in\mathcal E_N\), define
\begin{equation}\label{eq:endpoint-prefix-posterior}
 w(c;e)=P_0\bigl(\mathcal D_N^+,\ \pi X_{T_N}=e
                  \mid X\text{ with marks begins with }c\bigr).
\end{equation}
This is a deterministic function of \(c\) and \(e\).
All active departures in \(a_0\) lie strictly below height \(k-M\),
whereas all departures in \(a\) have height at least \(k\).
Lemma~\ref{lem:separated-weights} therefore gives
\[
 W(a_0a)=W(a_0)W(a).
\]
Consequently the raw probability of the prefix \(a_0a\) and the
endpoint event in \eqref{eq:endpoint-prefix-posterior} is
\(W(a_0)W(a)w(a_0a;e)\).

For a raw negative path starting at \((N-1,e)\), let \(D^-\) denote
the event that all its heights are at most \(N-1\).
Let \(\mathcal V(a)\) be the event that this path satisfies \(D^-\),
avoids the closed \(R\)-neighborhood of \(\Dep(a)\) for all time,
and visits the closed \(R\)-neighborhood of \(y(a)\).
Thus \(\mathcal V(a)\) asserts that the first contact along the word
\(a\) occurs at its terminal site.  If that site was an earlier
departure of \(a\), the event is empty.
The raw bridge interpretation and the conditioning by \(D^-\) give
the exact identity
\begin{equation}\label{eq:first-contact-prefix-sum}
 \begin{split}
 P(\mathcal C_{k,r,N})
 =\frac1{\bar u_N^+p_-}
 \sum_{a_0}W(a_0)
 \sum_{e\in\mathcal E_N}\sum_a
 W(a)w(a_0a;e)P_{(N-1,e)}(\mathcal V(a)).
 \end{split}
\end{equation}
Indeed, a pair of paths in \(\mathcal C_{k,r,N}\) determines a unique
\(a_0\) and a unique prefix \(a\) ending at its first contact after
\(T_k\).  Conversely every term represents exactly these conditions.
All word families are countable, and every summand is nonnegative.

\emph{A shared environment for the unconditioned prefixes.}
Fix \(a_0\) for the moment.  For a word \(a\), put
\[
 B(a)=\{x\in\ZZ^d:\dist_\infty(x,\Dep(a))\le R\}.
\]
The quenched probability of \(\mathcal V(a)\) can be computed from
the negative walk killed on arrival in \(B(a)\).  It requires eternal
survival, the upper-height restriction, and a visit near \(y(a)\).
It therefore uses only rows in \(B(a)^c\), including when these
requirements concern the entire infinite path.  Those rows are more
than distance \(R\) from every active departure of \(a\).
Finite-range independence gives
\begin{equation}\label{eq:first-contact-transfer}
 W(a)P_{(N-1,e)}(\mathcal V(a))
 =E_Q\left[p_\omega(a)
           P_{(N-1,e),\omega}(\mathcal V(a))\right].
\end{equation}

The right-hand side has the following interpretation.  Sample a fresh
environment with law \(Q\).  In it, run a raw comparison path
\(\widetilde X\) from \(z\) and a raw negative path
\(\widetilde\beta^{\,e}\) from \((N-1,e)\), independently
conditional on the environment.  Denote its law and expectation by
\(P_e^{\mathrm c},E_e^{\mathrm c}\), and write
\(P^{\mathrm c},E^{\mathrm c}\) for the common
\((\omega,\widetilde X)\) marginal, which does not depend on \(e\).
For fixed \(a_0,e\), the inner sum in
\eqref{eq:first-contact-prefix-sum} becomes
\begin{equation}\label{eq:transferred-prefix-sum}
 E_e^{\mathrm c}\left[
  \sum_a w(a_0a;e)
  \one_{\{\widetilde X\text{ begins with }a\}}
  \one_{\{\widetilde\beta^{\,e}\in\mathcal V(a)\}}
 \right].
\end{equation}
For every realized pair, at most one summand is nonzero: its terminal
site must be the first site of \(\widetilde X\) within distance
\(R\) of the full range of \(\widetilde\beta^{\,e}\).
This uniqueness is a pathwise fact and uses no stopping-time property
of the first contact.

\emph{The stopped posterior law.}
The factors \(w(a_0a;e)\) in
\eqref{eq:transferred-prefix-sum} still come from the original raw
positive law conditioned on \(a_0\).  In that law, let
\(\mathcal G_i\) record the marked continuation for \(i\) steps,
and form the martingales
\[
 w_i(e)=P_0(\mathcal D_N^+,\ \pi X_{T_N}=e
                  \mid a_0,\mathcal G_i),
 \qquad e\in\mathcal E_N.
\]
They are nonnegative and have sum at most one.  Stop the continuation
on its first hit of height \(k-1\) or \(N\).  Every departure before
that stop has height at least \(k\), so it is separated from the
active departures of \(a_0\).  The stopping arrival uses no row.
By factoring each finite stopped-prefix probability, the conditional
stopped continuation therefore has exactly the same path law as
\(\widetilde X\) stopped at
\[
 \zeta=T_{k-1}(\widetilde X)\wedge T_N(\widetilde X).
\]
Both stops are almost surely finite by uniform ellipticity and exit
from a slab of bounded height.

Evaluate the same deterministic functions
\eqref{eq:endpoint-prefix-posterior} along the comparison path until
\(\zeta\), and freeze them afterwards.  More explicitly, if
\(\widetilde X_{[0,j]}\) denotes its marked prefix of length \(j\),
set
\[
 \widetilde w_i(e)
 =w\bigl(a_0\widetilde X_{[0,i\wedge\zeta]};e\bigr).
\]
Equality of the stopped path laws shows that this vector has the
stopped martingale law just described.  Its values remain the original
conditional probabilities given \(a_0\); they are not conditional
probabilities given the fresh environment.  In particular, the value
at a hit of \(k-1\) retains its original meaning, and no identification
of continuation laws beyond that hit is needed.

Define
\[
 M_*(e)=\max_{0\le i\le\zeta}\widetilde w_i(e),
 \qquad S_*=\sum_{e\in\mathcal E_N}M_*(e).
\]
All endpoint coordinates are functions of the same stopped comparison
path.  We can therefore estimate their sum \(S_*\) on an event while
retaining only a logarithmic dependence on the number of endpoints.
Applying Lemma~\ref{lem:posterior-max} and using
\(\log(|\mathcal E_N|+1)\le C\log(2N)\) gives, for any event \(V\)
depending on the comparison path and its environment,
\begin{equation}\label{eq:stopped-posterior-event}
 E^{\mathrm c}[S_*\one_V]
 \le C\log(2N)\,P^{\mathrm c}(V)
                  \log\frac{\mathrm e}{P^{\mathrm c}(V)}.
\end{equation}
The constant is uniform over \(a_0\).  We have thus controlled the
endpoint factors while preserving a fresh shared environment in which
to estimate contacts.

\emph{Quenched exit probabilities.}
For an interior site \(y\), meaning \(k\le y_1<N\), write
\[
 q_y^\omega=P_{y,\omega}(T_{k-1}<T_N).
\]
Ellipticity gives \(0<q_y^\omega<1\).  For
\(t=2^{-j}\), \(j\ge1\), let \(V_t\) be the event that
\(\widetilde X\), at or after its first hit of height \(r\) and
before \(\zeta\), visits a site with \(q_y^\omega\ge t\).
With the environment fixed, stop at the first such visit.  The
conditional chance of exiting at \(k-1\) is then at least \(t\).
After averaging, stationarity and almost sure slab exit give
\begin{equation}\label{eq:first-contact-large-q}
 \begin{split}
 tP^{\mathrm c}(V_t)
 &\le P_z(T_r<T_{k-1}<T_N)\\
 &=u_{r-k}^+-u_{N-k}^+=\Delta.
 \end{split}
\end{equation}

Suppose a retained prefix ends at \(y\) with
\(q_y^\omega\in[t,2t)\).  Its endpoint lies at height at least
\(r\), so \(V_t\) occurs.  The negative path visits a site \(y'\)
within distance \(R\) of \(y\).  Because \(r-k\ge R\) and the
negative path satisfies \(D^-\), this site is also interior:
\(k\le y'_1\le N-1\).  A coordinate path from \(y\) to \(y'\)
of length at most \(dR\) stays inside the slab.  Following this path
and then exiting through the lower boundary shows that
\[
 q_y^\omega\ge\kappa^{dR}q_{y'}^\omega,
 \qquad q_{y'}^\omega\le c_Rt,
 \qquad c_R=2\kappa^{-dR}.
\]

For each fixed environment, the probability that a raw path from
\((N-1,e)\) satisfies \(D^-\) and visits any interior site with
\(q_{y'}^\omega\le c_Rt\) is at most \(c_Rt\), uniformly in \(e\).
Indeed, stop on the first visit to that set.  To satisfy \(D^-\)
after this visit, the path must next leave the slab through \(k-1\),
except on the null event of never exiting it.  The strong Markov property
bounds the conditional probability of this next exit by
\(q_{y'}^\omega\le c_Rt\) at the first-visit site, regardless of
earlier visits below \(k-1\).

\emph{Summing the exit-probability bins.}
We classify the transferred sum
\eqref{eq:transferred-prefix-sum} according to
\(q_{y(a)}^\omega\in[t,2t)\).  This classification is made in the
shared environment, after \eqref{eq:first-contact-transfer} has been
applied.  A retained prefix ends before \(\zeta\), so its posterior
factor is at most \(M_*(e)\).  By uniqueness of the retained prefix,
its contribution in this bin is pointwise bounded by \(M_*(e)\)
times the indicators of \(V_t\) and of the negative-path event in
the preceding paragraph.  Conditional on the environment, that
negative path is independent of \(\widetilde X\), and its event has
probability at most \(c_Rt\).  Summing over \(e\), whose comparison
path and environment have the same marginal, bounds this bin's total
contribution for fixed \(a_0\) by
\begin{equation}\label{eq:first-contact-bin}
 Ct\,E^{\mathrm c}[S_*\one_{V_t}].
\end{equation}

If \(\Delta=0\), \eqref{eq:first-contact-large-q} makes every term
zero.  Suppose \(\Delta>0\).  For \(t<\Delta\), use
\(E^{\mathrm c}S_*\le C\log(2N)\) and sum the geometric series in
\(t\).  The contribution is at most \(C\log(2N)\Delta\).
For \(t\ge\Delta\), combine
\eqref{eq:stopped-posterior-event} and
\eqref{eq:first-contact-large-q}.  Since
\(p\mapsto p\log(\mathrm e/p)\) is increasing on \([0,1]\), each
such bin contributes at most
\[
 C\log(2N)\,\Delta\log(\mathrm e t/\Delta).
\]
There are \(O(\log(\mathrm e/\Delta))\) bins in this second range,
so all bins together contribute at most
\(C\log(2N)\Delta\log^2(\mathrm e/\Delta)\).
Finally, the first-hit prefixes \(a_0\) are disjoint, hence
\(\sum_{a_0}W(a_0)\le1\), and
\eqref{eq:bridge-renewal} gives
\(\bar u_N^+p_-\ge\xi p_+p_->0\).
Substitution in \eqref{eq:first-contact-prefix-sum} proves
\eqref{eq:first-contact-bound}.
\end{proof}
